\documentclass[a4paper, 20pt,reqno]{article}
\usepackage[left=2cm,right=2cm,top=2cm,bottom=2cm]{geometry}
\usepackage[T2A]{fontenc}
% \DeclareMathSizes{10}{10}{10}{10}

\usepackage[russian]{babel}
\usepackage{soul}
\usepackage{gensymb}
\usepackage{amsfonts,amsmath,amssymb}
\usepackage{mathrsfs}
\usepackage{graphicx}
\usepackage[normalem]{ulem}
\usepackage[document]{ragged2e}
\usepackage{stmaryrd}
\usepackage{wrapfig}
\usepackage{fancyhdr}
\usepackage{floatflt}
\usepackage{python}
\usepackage{float}
\usepackage{amssymb}
\usepackage[most]{tcolorbox}
\usepackage{indentfirst}
\usepackage{setspace}
\usepackage{scrextend}
\usepackage{listings}
\usepackage{makecell,tabularx}
\usepackage{hyperref}
\usepackage{xcolor}

\newcommand{\mycopyright}{pluttan}
\newcommand{\docopyright}{$\mathfrak{Copyright}\ \mathfrak{\mycopyright} \logo$}
\newcommand{\rub}{{\rm{Р}\kern-.635em\rule[.5ex]{.52em}{.04em}\kern.11em}}

\definecolor{linkcolor}{HTML}{000000} 
\definecolor{urlcolor}{HTML}{0000FF} 

\hypersetup{pdfstartview=FitH,  linkcolor=linkcolor,urlcolor=urlcolor, colorlinks=true}

\definecolor{grey}{RGB}{40, 40, 40}

\renewcommand{\href}[1]{\url{#1}}

\definecolor{codegreen}{rgb}{0,0.6,0}
\definecolor{codegray}{rgb}{0.5,0.5,0.5}
\definecolor{codepurple}{rgb}{0.7,0,0.82}
\definecolor{mermaid}{rgb}{0.5,0.3,1}
\definecolor{backcolour}{rgb}{0.95,0.95,0.92}

\lstdefinestyle{mystyle}{
    backgroundcolor=\color{backcolour},   
    commentstyle=\color{codegreen},
    keywordstyle=\color{mermaid},
    numberstyle=\tiny\color{codegray},
    stringstyle=\color{codepurple},
    basicstyle=\ttfamily\footnotesize,
    breakatwhitespace=false,         
    breaklines=true,                 
    captionpos=b,                    
    keepspaces=true,                 
    numbers=left,                    
    numbersep=5pt,                  
    showspaces=false,                
    showstringspaces=false,
    showtabs=false,                  
    tabsize=2,
  extendedchars=true , % включаем не латиницу 
  escapechar=|, % |«выпадаем» в LATEX|
  frame=tb , % рамка сверху и снизу 
  commentstyle=\itshape , % шрифт для комментариев 
  stringstyle=\bfseries
}

\lstset{style=mystyle}

\lstdefinestyle{CommentStyle}{
    language=XML,
    %numbers=left, numberstyle=\tiny, stepnumber=1, numbersep=5pt,
    commentstyle=\color{red},
	basicstyle=\footnotesize\ttfamily,
	language={[ANSI]C++},
	keywordstyle=\bfseries,
	showstringspaces=false,
	morekeywords={include, printf},
	commentstyle={},
	escapeinside=§§,
	escapebegin=\begin{russian}\commentfont,
	escapeend=\end{russian},
    keywordstyle=\color{blue}\bfseries,
    morekeywords={align,begin},
    extendedchars=\true,
    tabsize=2
}
\lstdefinestyle{myLatexStyle}{
    language=c++,
    %backgroundcolor=\color{grey},
    numbers=left, numberstyle=\tiny, stepnumber=1, numbersep=5pt,
    commentstyle=\color{red},
    keywordstyle=\color{blue}\bfseries,
    morekeywords={align,begin},
    extendedchars=\true,
    tabsize=2
}

\lstdefinestyle{pmyLatexStyle}{
    language=java,
    %backgroundcolor=\color{grey},
    numbers=left, numberstyle=\tiny, stepnumber=1, numbersep=5pt,
    commentstyle=\color{red},
    keywordstyle=\color{blue}\bfseries,
    morekeywords={align,begin},
    extendedchars=\true,
    tabsize=2
}

\definecolor{block-gray}{gray}{0.90} % уровень прозрачности (1 - максимум)
\definecolor{yellow}{HTML}{F0FFFF}
\newtcolorbox{myquote}{colback=block-gray,grow to right by=-10mm,grow to left by=-10mm,boxrule=0pt,boxsep=0pt,breakable} % настройки области с изменённым фоном
\newtcolorbox{myquote2}{colback=yellow,grow to right by=-10mm,grow to left by=-10mm,boxrule=0pt,boxsep=0pt,breakable} % настройки области с изменённым фоном

\setlength{\parindent}{12,5mm}

\newcommand{\logo}{\vcenter{\hbox{\includegraphics[width=.8em]{/Users/pluttan/Documents/bw2.png}}}}
\onehalfspacing

\pagestyle{fancy}
\renewcommand{\sectionmark}[1]{\markright{#1}}
\fancyhf{} 
\fancyhead[R]{\bfseries\thepage}
\fancyhead[LO]{\docopyright}

\newcommand{\image}[2]{
	\begin{figure}[H]
		\center{\includegraphics[height=#2pt]{../img/#1} }
    \end{figure}
}

\newcommand{\dotitle}[2]{

\thispagestyle{empty}
\sloppy{
  \scriptsize{
    \line(6,0){0}

    \centering \docopyright

    \centering Привет! Меня зовут Андрей, я создаю свою ботву, этот файл малая ее часть. 
    
    \centering Пользоваться и распространять файлы конечно же можно. Если вы нашли ошибку в файле, можете 
    
    \centering исправить ее в исходном коде и подать на слияние или просто написать в issue. 

    \centering {\bf{Так же вы можете купить распечатанную версию данного файла в виде книжки.}
    
    \centering По всем вопросам писать в ВК.}
    
    \centering Приятного бота)

    \line(6,0){300}

    \centering GitHub: \href{https://github.com/pluttan}

    \centering VK: \href{https://vk.com/pluttan}

    \line(6,0){0}
  }}

\begin{figure}[H]
		\center{\includegraphics[height=100pt]{/Users/pluttan/Documents/forMyDocs/qr2.png} }
\end{figure}

\topskip=-200pt
\vspace*{130pt}
\begin{Huge}
  \textbf{
    \begin{center}
        #1
    \end{center}
  }
  {\begin{center} 
        #2
    \end{center}}
\end{Huge}
\vspace*{300pt}
  \begin{flushright}
    Над файлом работали: \\
    \mycopyright
  \end{flushright}
\newpage
\pagestyle{fancy}
\renewcommand{\sectionmark}[1]{\markright{#1}}
\fancyhf{} 
\fancyhead[R]{\bfseries\thepage}
\fancyhead[LO]{\docopyright}
}
\renewcommand{\sectionmark}[1]{\markright{#1}}
\renewcommand{\sectionmark}[1]{\markright{#1}}
% \makeatletter
%      \renewcommand*\l@section{\@dottedtocline{1}{0em}{1.8em}}
%      \renewcommand*\l@subsection{\@dottedtocline{2}{1.5em}{2.0em}}
%      \renewcommand*\l@subsubsection{\@dottedtocline{3}{4.3em}{3.0em}}
% \makeatother
\newcommand{\toc}{
\newpage
\renewcommand{\contentsname}{Оглавление}
\large{\tableofcontents}
\newpage
\large
}


% It's XeLaTeX, if you can't compile comment this line 
\usepackage{fontspec}\setmainfont{Times New Roman} 


\begin{document}

\dotitle{Подготовка к РК1}{Математический анализ}{}

\section{Определения}
    \subsection{Сформулируйте определение окрестности точки $x \in \mathbb{R}$.}
    111
    Окрестностью $U(x_0)$ точки $x_0$ называют любой интервал, содержащий эту точку

    \subsection{Сформулируйте определение $\epsilon$-окрестности точки $x \in \mathbb{R}$. }

    $\epsilon$-окрестностью $U_\epsilon(x_0)$ точки $x_0$ называют интервал с центром в 
    $x_0$ и длиной $2\epsilon$ т.е. $$U_\epsilon(x_0) = (x_0 - \epsilon, x_0 + \epsilon) = 
    \{x \in \mathbb{R},  \epsilon>0 : x_0 - \epsilon < x < x_0 + \epsilon\}$$.

    \subsection{Сформулируйте определение окрестности $+\infty$. }

    Окрестностью точки $+\infty$ называют интервал вида $(a, +\infty)$, где $a$ — произвольное действительное число.
    $$U(+\infty) = \{x \in \mathbb{R}: x > a\}, a > 0$$

    \subsection{Сформулируйте определение окрестности $-\infty$. }

    Окрестностью точки $-\infty$ называют интервал вида $(-\infty, a)$, где $a$ — произвольное действительное число.
    $$U(-\infty) = \{x \in \mathbb{R}: x < -a\}, a > 0$$

    \subsection{Сформулируйте определение окрестности $\infty$. }

    Окрестностью бесконечности $\infty$ «без знака»  называют объединение двух бесконечных интервалов 
    $(-\infty, -a) \cup (a, +\infty)$, где a — произвольное действительное число.
    $$U(\infty) = \{x \in \mathbb{R}: |x| > a\}, a > 0$$

    \subsection{Сформулируйте определение предела последовательности.}
    
    Число a называется пределом последовательности $\{x_n\}, n \to +\infty$, если для любого 
    сколь угодно малого $\epsilon > 0$ существует номер $N = N(\epsilon)$ такой, что 
    если порядковый номер члена последовательности $n  \geqslant  N$, то имеет место неравенство 
    $|x_n - a| < \epsilon$. 
    $$\lim\limits_{n \to \infty} x_n = a \iff \ \forall \epsilon > 0 \  \exists N = N(\epsilon) \in 
    \mathbb{N} : \forall n > N(\epsilon) \longrightarrow |x_n - a| < \epsilon$$
    
    \subsection{Сформулируйте определение сходящейся последовательности.}

    Последовательность, придел которой существует и конечен при $n \to \infty$. Поскольку 
    неравенство $|a - x_n| < \epsilon$ эквивалентно неравенству 
    $a - \epsilon < x_n < a + \epsilon$, то все элементы сходящейся последовательности 
    за исключением конечного их числа при любом $\epsilon > 0$ лежат в $\epsilon$-окрестности 
    точки $a$.

    \subsection{Сформулируйте определение ограниченной последовательности.}

    Последовательность $\{x_n\}$ называется ограниченной снизу, если существует число $c_1$
    такое, что $x_n  \geqslant  c_1$ при всех $n = 1, 2, ...$

    Последовательность $\{x_n\}$ называется ограниченной сверху, если существует число $c_2$
    такое, что $x_n  \leqslant  c_2$ при всех $n = 1, 2, ...$

    Последовательность $\{x_n\}$ ограниченная как сверху, так и снизу, называется ограниченной,
    то есть $c_1  \leqslant  x_n  \leqslant  c_2$ при всех $n = 1, 2, ...$

    $$\exists M > 0, M \in \mathbb{R}: \forall n \in \mathbb{N} \longrightarrow |x_n|  \leqslant  M 
    (m  \leqslant  x_n  \leqslant  M, m \in \mathbb{R})$$

    \subsection{Сформулируйте определение монотонной последовательности.}
    
    Последовательность называется монотонной, если она неубывающая($x_{n+1} \geqslant x_n$), возрастающая($x_{n+1}>x_n$), 
    невозрастающая($x_{n+1} \leqslant x_n$) или убывающая($x_{n+1}<x_n$) для $\forall n \in \mathbb{N}$

    \subsection{Сформулируйте определение возрастающей последовательности.}

    Последовательность называется возрастающей, если $x_{n+1} > x_n \forall n \in \mathbb{N}$
        
    \subsection{Сформулируйте определение убывающей последовательности.}

    Последовательность называется убывающей, если $x_{n+1} < x_n \forall n \in \mathbb{N}$
        
    \subsection{Сформулируйте определение невозрастающей последовательности.}
    
    Последовательность называется невозрастающей, если $x_{n+1} \leqslant x_n \forall n \in \mathbb{N}$
        
    \subsection{Сформулируйте определение неубывающей последовательности.}
    
    Последовательность называется неубывающей, если $x_{n+1} \geqslant x_n \forall n \in \mathbb{N}$
    
    \subsection{Сформулируйте определение фундаментальной последовательности.}
    
    Последовательность $\{x_n\}$ называется фундаментальной, если для любого $\epsilon > 0$ 
    существует номер $N = N(\epsilon)$ такой, что для всех $m  \geqslant  N$ и $n  \geqslant  N$ выполняется 
    неравенство $| x_m - x_n | < \epsilon$.

    \subsection{Сформулируйте критерий Коши существования предела последовательности.}

    Для того, чтобы последовательность была сходящейся, необходимо и достаточно, чтобы она была фундаментальной.

    \subsection{Сформулируйте определение по Гейне предела функции.}

    Пусть функция $f(x)$ определена в проколотой окрестности $\mathring U(x_0)$ точки $x_0$. 
    Число $a$ называется пределом функции $f(x)$ при $x \to x_0$, если для любой последовательности 
    $\{x_n\}$ точек из $\mathring U(x_0)$ , для которой $\lim\limits_{n \to \infty} x_n = x_0$,
    выполняется равенство $\lim\limits_{n \to \infty} f(x_n) = a$.

    \subsection{Сформулируйте определение бесконечно малой функции.}

    Функция $f(x)$ называется бесконечно малой при $x \to x_0, x_0 \in \mathbb{R}$, если 
    $\lim\limits_{x \to x_0} f(x) = 0$.

    \subsection{Сформулируйте определение бесконечно большой функции.}
    
    Функция $f(x)$ называется бесконечно большой при $x \to x_0, x_0 \in \mathbb{R}$, 
    если $\lim\limits_{x \to x_0} f(x) = +\infty$.

    \subsection{Сформулируйте определение бесконечно малых функций одного порядка.}

    Если существует конечный отличный от нуля предел $\lim\limits_{x \to x_0}\frac{f(x)}{g(x)} = C \ne 0$,
    то говорят, что $f(x)$ и $g(x)$ являются при $x \to x_0$ бесконечно малыми одного порядка и пишут $f(x) = O(g(x))$.

    \subsection{Сформулируйте определение несравнимых бесконечно малых функций.}

    Если при $x \to x_0$ не существует ни конечного, ни бесконечного предела 
    отношения $\frac{f(x)}{g(x)}$ или $\frac{g(x)}{f(x)}$, то говорят, что $f(x)$ и $g(x)$ не сравнимы при $x \to x_0$.

    \subsection{Сформулируйте определение эквивалентных бесконечно малых функций.}
    
    В случае $C = 1$, т.е. если $\lim\limits_{x \to x_0}\frac{f(x)}{g(x)} = 1$, 
    функции $f(x)$ и $g(x)$ называют эквивалентными бесконечно малыми и пишут $f(x) \sim g(x)$, при $x \to x_0$.

    \subsection{Сформулируйте определение порядка малости одной функции относительно другой.}

    Пусть $f(x)$ и $g(x)$ бесконечно малые при $x \to x_0$. Если при некотором $k$
    бесконечно малые $f(x)$ и $(g(x))^k$ являются бесконечно малыми одного порядка,
    то говорят, что $f(x)$ имеет порядок малости $k$ по сравнению с $g(x)$ при $x \to x_0$.

    \subsection{Сформулируйте определение приращения функции.}

    Приращением функции называют $\Delta f(x_0) = f(x) - f(x_0) = f(x_0 + \Delta x) - f(x_0)$.

    \subsection{Сформулируйте определение непрерывности функции в точке (любое).}

    Функция $f(x)$ называется непрерывной в точке $x_0$, если выполняется 3 пункта:
    \begin{enumerate}
        \item $\exists\lim\limits_{x \to x_0 + 0}f(x),\lim\limits_{x \to x_0 - 0}f(x)$
        \item $\exists f(x_0)$
        \item $f(x_0) = \lim\limits_{x \to x_0 + 0}f(x) = \lim\limits_{x \to x_0 - 0}f(x)$
    \end{enumerate} 

    \subsection{Сформулируйте определение непрерывности функции на интервале.}

    Функция непрерывна на интервале, если она непрерывна в каждой точке этого интервала.

    \subsection{Сформулируйте определение непрерывности функции на отрезке.}

    Функция непрерывна на отрезке $[a,b]$, если она непрерывна на интервале $(a, b)$,
    в точке $x = a$ слева, и в точке $x = b$ справа.

    \subsection{Сформулируйте определение точки разрыва.}

    Точка $x_0$ называется точкой разрыва функции $f(x)$, если не выполняется хотя бы один из 3 пунктов:
    \begin{enumerate}
        \item $\exists\lim\limits_{x \to x_0 + 0}f(x),\lim\limits_{x \to x_0 - 0}f(x)$
        \item $\exists f(x_0)$
        \item $f(x_0) = \lim\limits_{x \to x_0 + 0}f(x) = \lim\limits_{x \to x_0 - 0}f(x)$
    \end{enumerate} 

    \subsection{Сформулируйте определение точки устранимого разрыва.}

    Если $f(x_0) \ne \lim\limits_{x \to x_0 + 0}f(x) = \lim\limits_{x \to x_0 - 0}f(x)$, то 
    $x_0$ называют точкой устранимого разрыва функции $f(x)$.

    \subsection{Сформулируйте определение точки разрыва I-го рода.}
    
    Если $x_0$ — точка разрыва функции $f(x)$, и существуют конечные пределы
    $\lim\limits_{x \to x_0-} f(x) = f(x_0 - 0)$ и $\lim\limits_{x \to x_0+} f(x) = f(x_0 + 0)$,
    то $x_0$ называется точкой разрыва первого рода.

    \subsection{Сформулируйте определение точки разрыва II-го рода.}
    
    Функция $f(x)$ имеет точку разрыва второго рода при $x=x_0$, если по крайней мере один 
    из односторонних пределов не существует или равен бесконечности.

\newpage
\section{Определение предела по Коши}
    \subsection{Сформулируйте определение по Коши $\lim\limits_{x \to 0} f (x) = b$, где $b \in \mathbb{R}$. Приведите 
    соответствующий пример (с геометрической иллюстрацией).}

    Пусть $f(x)$ определена в проколотой окрестности $\mathring U(x_0)$ точки $x_0 = 0$. Число $b \in \mathbb{R}$ 
    называется приделом функции $f(x)$ в точке $x \to x_0$, если для любого $\epsilon > 0$ существует число 
    $\delta = \delta(\epsilon) > 0$ такое, что  для всех $x \in \mathring U_\delta(x_0)$ выполняется неравенство 
    $|f(x) - b| < \epsilon$
    \image{1.png}{200}

    \subsection{Сформулируйте определение по Коши $\lim\limits_{x \to a} f (x) = +\infty$, где $a \in \mathbb{R}$. Приведите 
    соответствующий пример (с геометрической иллюстрацией).}

    Пусть $f(x)$ определена в проколотой окресности $\mathring U(x_0), x_0 = a$. Тогда 
    $\lim\limits_{x \to x_0} f(x) = +\infty$, если существует сколь угодно малое $\epsilon > 0$, 
    такое что найдется число $\delta = \delta(\epsilon) > 0$, при котором если $x \in \mathring U_\delta(x_0) \ 
    0 < |x - x_0| < \delta$, то справедливо неравенство $f(x) > \epsilon$ 
    \image{2.png}{200}

    \subsection{Сформулируйте определение по Коши $\lim\limits_{x \to \infty} f (x) = 0$. Приведите соответствующии 
    пример (с геометрической иллюстрацией).}

    Пусть $f(x)$ определена в окресности $+\infty$. Тогда число 0 является пределом функции
    $f(x)$, если для любого сколь угодно малого $\epsilon > 0$ найдется такое $\delta = \delta(\epsilon) > 0$, что
    если $\forall x > \delta$, то верно неравенство $|f(x)| < \epsilon$
    \image{3.png}{200}

    \subsection{Сформулируйте определение по Коши $\lim\limits_{x \to a-0} f (x) = -\infty$, где $a \in 
    \mathbb{R}$. Приведите соответствующий пример (с геометрической иллюстрацией).}

    $\lim\limits_{x \to a-0} f (x) = -\infty$ - левосторонний предел функции $f(x)$ при $x \to a$, если 
    для любого сколь угодно малого $\epsilon > 0$ найдется такое $\delta = \delta(\epsilon) > 0$, что если
    $x \in \mathring U_\delta^-(x_0)$, то верно неравенство $f(x) < -\epsilon$
    \image{4.png}{200}

\newpage
\section{Теоремы}

    \subsection{Сформулируйте теорему об ограниченности сходящейся числовой последовательности. }

    Всякая сходящаяся последовательность ограничена.
    $$\exists \lim\limits_{x \to \infty} x_n = a, a \in \mathbb{R} \rightarrow 
    \exists M > 0: \forall n \in \mathbb{N} \  x_n \in M$$

    \subsection{Сформулируйте теорему о связи функции, ее предела и бесконечно малой. }

    Если существует конечный $\lim\limits_{x \to a} f(x) = A$, то $f(x)$ представлена в виде
    $f(x) = A + \alpha  (x)$, где $\lim\limits_{x \to a} \alpha  (x) = 0$ - бесконечно малая при $x \to a$

    \subsection{Сформулируйте теорему о сумме конечного числа бесконечно малых функций. }

    Сумма конечного числа функций, являющихся бесконечно малыми, при $x \to a$, 
    есть величина бесконечно малая при $x \to a$
    $$\exists \lim\limits_{x \to a} \alpha _k (x) = 0 \rightarrow \lim\limits_{x \to a} 
    \sum\limits_{k=1}^n \alpha _k (x) = 0$$


    \subsection{Сформулируйте теорему о произведении бесконечно малой на ограниченную функцию. }
    
    Если функция $\alpha  (x)$ бесконечно малая при $x \to a$, а $f(x)$ ограниченная функция, то
    $\alpha  (x) \times f(x)$ бесконечно малая, при $x \to a$.

    \subsection{Сформулируйте теорему о связи бесконечно малой и бесконечно большой функций. }

    Если $f(x)$ бесконечно большая при $x \to a$, то $\frac{1}{f(x)}$ бесконечно малая при $x \to a$.
    Если $\alpha  (x)$ бесконечно малая при $x \to a$ и отличная от нуля, то $\frac{1}{\alpha (x)}$ 
    бесконечно большая при $x \to a$

    \subsection{Сформулируйте теорему о необходимом и достаточном условии эквивалентности бесконечно малых. }

    Две бесконечно малые ф-ции при $x \to a$ эквивалентны $\iff$ их разность есть бесконечно
    малая более высокого порядка, чем каждая из них

    $$(f(x) \sim g(x) \ x \to a) \iff (f(x) - g(x) = o(f(x))) \vee (f(x) - g(x) = o(g(x)))$$


    \subsection{Сформулируйте теорему о сумме бесконечно малых разных порядков. }
    
    Если функции $\alpha  (x), \beta (x), ..., \gamma (x)$ бесконечно малые при $x \to a$,
    то $\alpha  (x) + \beta (x) + ... + \gamma (x) \sim \alpha  (x)$ при $x \to a$,
    где $\lim\limits_{x \to a} \frac{\beta (x)}{\alpha (x)} = 0; ...
    \lim\limits_{x \to a} \frac{\gamma (x)}{\alpha (x)} = 0$
    
\end{document}